\section{Scope and Evidence}

\subsection{Evidence classes}
\noindent
\textbf{Analytical evidence.}
The analytical record includes wavelength, ideal half-wave length, the thin resonant
half-wave dipole pattern, numerical directivity integration, and analytical HPBW.

\noindent
\textbf{Simulated evidence.}
The simulation record includes CST geometry screenshots, three length-tuning
$S_{11}$ plots, and two-dimensional and three-dimensional far-field figures. The final
$S_{11}$ figure contains an exact marker at \SI{0.99648}{GHz}.

\noindent
\textbf{Measured evidence.}
No prototype, VNA calibration, balun, feed cable, chamber range, gain-transfer test,
or uncertainty budget is included. No measured claim is made.

\subsection{Public-edition boundaries}
Administrative course information and student identifiers were removed. Native CST
binary models were not present in the supplied archive and are therefore not included.
The public package preserves the supplied figures, an independent calculation script,
a run manifest, and explicit numerical limitations.

\subsection{Claim discipline}
The report distinguishes among three levels of numerical confidence:
\begin{enumerate}[leftmargin=2em]
\item \textbf{Exact exported marker:} final resonance and final minimum $S_{11}$.
\item \textbf{Displayed figure value:} far-field directivity and angular-width annotations.
\item \textbf{Plot-derived estimate:} resonance and minimum depth for the two preliminary lengths.
\end{enumerate}
This distinction prevents screenshot estimates from being presented with unsupported
precision.
